%!TEX ROOT = thesis.tex

% =========================================================
% CHAPTER 5: RESEARCH DESIGN AND METHODOLOGY
% =========================================================
\chapter{RESEARCH DESIGN AND METHODOLOGY}
\label{chap:methodology}

This chapter establishes the rigorous methodological framework governing the development, optimization, and evaluation of the proposed Nano-Mamba U-Net. Moving beyond empirical heuristics, this section structurally justifies every engineering decision—from addressing the inherent physical artifacts of Magnetic Resonance Imaging (MRI) during preprocessing to formulating a mathematically robust optimization objective capable of mitigating severe anatomical class imbalances. Furthermore, this chapter provides an exhaustive architectural and mathematical deconstruction of the deep learning system itself, explicitly fulfilling the "Research Design" mandate by detailing the continuous-to-discrete state space transformations that empower the model.

\section{Dataset Justification and Clinical Heterogeneity}
The empirical validation of any deep learning architecture in medical imaging is profoundly dependent on the quality, diversity, and clinical realism of the utilized dataset. Deep neural networks are notoriously prone to overfitting to the specific scanner characteristics or healthy anatomical priors of a limited cohort \cite{campello2021multi}. To ensure the proposed model's robustness and clinical translatability, this research exclusively leverages the Automated Cardiac Diagnosis Challenge (ACDC) dataset \cite{bernard2018deep}.

The ACDC cohort consists of 100 distinct patients whose Cardiac Magnetic Resonance (CMR) cine sequences were acquired using 1.5T and 3.0T Siemens scanners. The defining characteristic of this dataset is its deliberate inclusion of severe pathological heterogeneity. The dataset is equally partitioned into five clinically distinct groups: Normal subjects (NOR), patients with previous Myocardial Infarction (MINF), Dilated Cardiomyopathy (DCM), Hypertrophic Cardiomyopathy (HCM), and Abnormal Right Ventricle (ARV). From an algorithmic perspective, these pathologies introduce extreme intra-class variability. For instance, in HCM patients, the myocardial wall is pathologically thickened, often blurring the boundaries with the papillary muscles, whereas in MINF patients, the presence of scar tissue severely alters the native magnetic resonance signal intensities and local image textures. The fundamental objective is to train the Nano-Mamba U-Net to learn highly abstract, pathology-invariant spatial representations to accurately segment the Right Ventricle (RV), Myocardium (MYO), and Left Ventricle (LV) across these diverse physical deformations.

\section{Data Harmonization and Spatial Resampling}
Medical imaging data generated in real-world clinical environments is inherently unstandardized. Variability in patient positioning, scanner FOV (Field of View), and institutional acquisition protocols lead to severe inconsistencies in physical voxel spacing and tensor dimensions. As demonstrated by the seminal work on the nnU-Net architecture \cite{isensee2021nnu}, the systematic harmonization of spatial metadata is arguably as critical to segmentation accuracy as the neural network architecture itself. 

The sequential progression of the data tensors through the preprocessing pipeline is structurally represented in Figure \ref{fig:preprocessing_pipeline}. This automated harmonization was rigorously implemented utilizing the MONAI (Medical Open Network for AI) framework \cite{cardoso2022monai}.

\begin{figure}[hbt!]
\centering
\begin{tikzpicture}[
    node distance=1.5cm and 2cm,
    block/.style={rectangle, draw=black!80, thick, fill=gray!10, text width=4.5cm, align=center, rounded corners=3pt, minimum height=1.2cm},
    arrow/.style={->, >=stealth, thick, draw=black!70}
]
% Nodes
\node[block] (raw) {\textbf{Raw NIfTI Volume}\\ Heterogeneous Spacing};
\node[block, right=of raw] (channel) {\textbf{Channel Expansion}\\ $\mathbb{R}^{H \times W \times D} \rightarrow \mathbb{R}^{1 \times H \times W \times D}$};
\node[block, below=of channel] (resample) {\textbf{Spatial Resampling}\\ Target: $256 \times 256 \times 16$};
\node[block, left=of resample] (intensity) {\textbf{Intensity Scaling}\\ Normalizing MR Physics};
\node[block, below=of intensity] (tensor) {\textbf{PyTorch Tensor}\\ Cuda/Hardware Ready};
% Arrows
\draw[arrow] (raw) -- (channel);
\draw[arrow] (channel) -- node[right, font=\footnotesize] {Continuous/Discrete} (resample);
\draw[arrow] (resample) -- node[above, font=\footnotesize] {Voxel Alignment} (intensity);
\draw[arrow] (intensity) -- node[right, font=\footnotesize] {Optimization} (tensor);
\end{tikzpicture}
\caption{The chronological data harmonization pipeline. Raw 3D MRI volumes undergo structural standardization to resolve spatial anisotropy and signal variability prior to network ingestion.}
\label{fig:preprocessing_pipeline}
\end{figure}

\subsection{Mathematical Formulation of Trilinear Interpolation}
To ensure structural compatibility with the Nano-Mamba tensor dimensions, all variable-sized volumes were spatially resampled to a standardized discrete grid of $256 \times 256 \times 16$. Because image intensities represent a continuous physiological signal captured at discrete spatial coordinates, mapping them to a new grid requires robust mathematical interpolation to prevent severe aliasing artifacts. 

For the continuous MRI intensity volumes, Trilinear Interpolation was applied. Let $V(x, y, z)$ denote the interpolated voxel intensity at a continuous Cartesian coordinate. By identifying the eight bounding discrete neighboring voxels $V_{000}$ through $V_{111}$ within the original grid, the continuous trilinear approximation is formulated as:
\begin{equation}
    V(x,y,z) \approx \sum_{i=0}^{1} \sum_{j=0}^{1} \sum_{k=0}^{1} V_{ijk} \cdot (1 - |x - x_i|) \cdot (1 - |y - y_j|) \cdot (1 - |z - z_k|),
\end{equation}
where $x_i, y_j, z_k$ represent the normalized spatial coordinates of the respective neighbors. 

Conversely, the ground-truth segmentation masks consist of strictly categorical integers (e.g., $1$ for RV, $2$ for MYO). Applying continuous polynomial interpolation to these masks would yield meaningless floating-point values, thereby destroying the categorical class definitions. Consequently, the label matrices are resampled utilizing Nearest-Neighbor Interpolation. This non-linear operation assigns the class of the closest original voxel by minimizing the Euclidean distance:
\begin{equation}
    V_{mask}(x, y, z) = V \left( \arg\min_{(x_i, y_i, z_i)} \sqrt{(x-x_i)^2 + (y-y_i)^2 + (z-z_i)^2} \right).
\end{equation}

\subsection{Addressing MR Intensity Heterogeneity}
Unlike Computed Tomography (CT), which relies on a standardized physical radiodensity scale (Hounsfield Units), the intensity values in MRI are relative and highly dependent on the specific acquisition parameters (e.g., echo time, repetition time) and the intrinsic magnetic field strength \cite{nyul1999on}. Consequently, the exact same myocardial tissue can manifest with drastically different numerical pixel intensities across different patients, creating severe instability during gradient descent optimization.

To construct a normalized feature space for the neural network, the spatial intensities are systematically bounded. The Min-Max scaling function $S(\mathcal{I})$ projects the continuous intensity range into a normalized subset $[0, 1]$:
\begin{equation}
    S(\mathcal{I}_{xyz}) = \frac{\mathcal{I}_{xyz} - \min(\mathcal{I})}{\max(\mathcal{I}) - \min(\mathcal{I}) + \epsilon},
\end{equation}
where $\epsilon = 10^{-8}$ serves as a negligible numerical stability constant designed to rigorously prevent division-by-zero singularities in the event of homogeneously black image slices.

\section{System Architecture Design: The Nano-Mamba U-Net}
The core "Research Design" of this thesis is the formulation of the Nano-Mamba architecture. The network adheres to a symmetric Encoder-Bottleneck-Decoder topology, fundamentally inspired by the standard 3D U-Net \cite{cicek20163d}, but critically deviates in its deepest representational layer. The architecture leverages consecutive 3D convolutional blocks in the contracting path to extract localized spatial hierarchies, and utilizes transposed convolutions in the expanding path for precise anatomical localization. The core innovation—the Spatio-Temporal Mamba Bottleneck—is strategically positioned at the apex of the contracting path, where the spatial dimensions are maximally compressed, and the semantic channel capacity is at its peak.

To visually articulate the complex forward propagation and information fusion mechanisms, Figure \ref{fig:full_architecture} presents the complete topological schematic of the proposed network.

\begin{figure}[hbt!]
\centering
\resizebox{0.95\textwidth}{!}{
\begin{tikzpicture}[
    node distance=1.2cm and 1.5cm,
    layer/.style={rectangle, draw=black!80, thick, align=center, rounded corners=2pt, minimum height=0.8cm},
    enc/.style={layer, fill=blue!10},
    dec/.style={layer, fill=green!10},
    mamba/.style={layer, fill=red!10, draw=red, very thick},
    arrow/.style={->, >=stealth, thick},
    skip/.style={->, >=stealth, thick, dashed, draw=blue}
]

% Encoder Path (Left Side)
\node[enc] (in) {\textbf{Input 3D MRI}\\ $1 \times 256 \times 256 \times 16$};
\node[enc, below=of in] (e1) {\textbf{Encoder 1}\\ Conv3D + ReLU\\ $16 \times 256 \times 256 \times 16$};
\node[enc, below=of e1] (e2) {\textbf{Encoder 2}\\ MaxPool $\downarrow$ + Conv3D\\ $32 \times 128 \times 128 \times 8$};
\node[enc, below=of e2] (e3) {\textbf{Encoder 3}\\ MaxPool $\downarrow$ + Conv3D\\ $64 \times 64 \times 64 \times 4$};

% Bottleneck (Bottom Center)
\node[mamba, below right=1cm and 1cm of e3] (bot) {\textbf{Spatio-Temporal Mamba Bottleneck (128 Channels)}\\ Tensor Serialization $\rightarrow$ Selective Scan $\rightarrow$ Volumetric Reshaping\\ $\mathbb{R}^{128 \times 32 \times 32 \times 2}$};

% Decoder Path (Right Side)
\node[dec, above right=1cm and 1cm of bot] (d3) {\textbf{Decoder 3}\\ UpConv $\uparrow$ + Concat + Conv3D\\ $64 \times 64 \times 64 \times 4$};
\node[dec, above=of d3] (d2) {\textbf{Decoder 2}\\ UpConv $\uparrow$ + Concat + Conv3D\\ $32 \times 128 \times 128 \times 8$};
\node[dec, above=of d2] (d1) {\textbf{Decoder 1}\\ UpConv $\uparrow$ + Concat + Conv3D\\ $16 \times 256 \times 256 \times 16$};
\node[dec, above=of d1] (out) {\textbf{Output Mask}\\ $1 \times 1 \times 1$ Conv3D\\ $4 \times 256 \times 256 \times 16$};

% Connections
\draw[arrow] (in) -- (e1);
\draw[arrow] (e1) -- (e2);
\draw[arrow] (e2) -- (e3);
\draw[arrow] (e3) |- (bot);
\draw[arrow] (bot) -| (d3);
\draw[arrow] (d3) -- (d2);
\draw[arrow] (d2) -- (d1);
\draw[arrow] (d1) -- (out);

% Skip Connections
\draw[skip] (e3) -- node[above, font=\scriptsize] {Skip Connection (Concat)} (d3);
\draw[skip] (e2) -- node[above, font=\scriptsize] {Skip Connection (Concat)} (d2);
\draw[skip] (e1) -- node[above, font=\scriptsize] {Skip Connection (Concat)} (d1);

\end{tikzpicture}
}
\caption{The complete architectural topology of the Nano-Mamba U-Net. Blue blocks represent the contracting spatial encoder, green blocks represent the expanding decoder, and the central red block denotes the linear-time Selective State Space Bottleneck. Dashed blue lines indicate spatial skip connections preserving high-frequency anatomical boundaries.}
\label{fig:full_architecture}
\end{figure}

\subsection{Micro-Architectural Components}
Each stage of the encoder and decoder pathways utilizes a "DoubleConv" block, comprising two consecutive $3 \times 3 \times 3$ convolutional layers. Mathematically, these layers are augmented with 3D Batch Normalization (BN) \cite{ioffe2015batch} and Rectified Linear Unit (ReLU) activation functions. 

The integration of 3D BN is critical for stabilizing the internal covariate shift when processing heterogeneous MRI volumes. For a mini-batch of feature maps $\mathcal{B}$, the normalization for a specific channel is formulated as:
\begin{equation}
    \hat{x} = \frac{x - E[\mathcal{B}]}{\sqrt{Var[\mathcal{B}] + \epsilon}},
\end{equation}
followed by a learnable affine transformation $y = \gamma \hat{x} + \beta$. This ensures that the gradient flow remains within a stable numerical range, preventing the vanishing gradient problem in deep 3D networks. Following normalization, the non-linearity is introduced via the ReLU function:
\begin{equation}
    f(x) = \max(0, x).
\end{equation}
By maintaining a constant gradient of 1 for all positive activations, ReLU facilitates sparse representations and accelerates the convergence of the Adam optimizer.

\section{The Spatio-Temporal Mamba Bottleneck}
The empirical superiority of the proposed architecture relies entirely on the successful mathematical integration of the Selective State Space Model \cite{gu2023mamba} into the highly compressed latent space.

\subsection{Volumetric Tensor Serialization}
Unlike localized 3D convolutions that naturally process grid-like spatial structures, the Mamba architecture fundamentally requires a 1D serialized sequence to perform its hardware-aware parallel scan. Let the highly compressed spatial feature map outputted by Encoder 3 and passed through the initial bottleneck convolution be denoted as a 5-dimensional continuous tensor:
\begin{equation}
    \mathcal{X} \in \mathbb{R}^{B \times C \times H \times W \times D},
\end{equation}
where $B$ represents the mini-batch dimension, $C=128$ is the optimal semantic channel depth (as definitively proven by our subsequent ablation studies), and $H, W, D$ are the spatial coordinates ($32 \times 32 \times 2$). 

To ingest this volumetric data into the State Space ODEs, the tensor is systematically flattened across its three spatial dimensions into a continuous 1D sequence of length $L$, where $L = H \times W \times D$. This mathematical transformation is expressed as:
\begin{equation}
    \mathbf{x}_{seq} = \text{Reshape}(\mathcal{X}) \in \mathbb{R}^{B \times C \times L}.
\end{equation}
To strictly align with the causal, autoregressive sequence processing requirements of the Mamba underlying CUDA kernels, the tensor must be transposed to treat the sequence length as the primary operational dimension:
\begin{equation}
    \mathbf{x}' = \text{Transpose}(\mathbf{x}_{seq}) \in \mathbb{R}^{B \times L \times C}.
\end{equation}
This serialization essentially forces the model to traverse the 3D cardiac volume as a temporal sequence, sequentially mapping the spatial voxels across the temporal cardiac cycle.

\subsection{Continuous-Time State Space ODEs and Discretization}
Once serialized, the sequence is processed by the State Space mechanism. Fundamentally, Mamba maps the 1D input sequence $x(t) \in \mathbb{R}$ to an output sequence $y(t) \in \mathbb{R}$ through an implicit, high-dimensional hidden latent state $h(t) \in \mathbb{R}^N$. This continuous-time dynamic mapping is governed by a set of linear Ordinary Differential Equations (ODEs):
\begin{equation}
    h'(t) = \mathbf{A}h(t) + \mathbf{B}x(t),
\end{equation}
\begin{equation}
    y(t) = \mathbf{C}h(t).
\end{equation}
Here, $\mathbf{A} \in \mathbb{R}^{N \times N}$ is the critically important state transition matrix that dictates the memory decay of the system. $\mathbf{B} \in \mathbb{R}^{N \times 1}$ and $\mathbf{C} \in \mathbb{R}^{1 \times N}$ are learnable projection matrices.

Because modern deep learning hardware operates on discrete, quantized tensors rather than continuous differential functions, the continuous ODEs must be analytically discretized. Mamba achieves this through the Zero-Order Hold (ZOH) assumption, introducing a hardware-aware temporal step size parameter $\mathbf{\Delta}$. Applying the ZOH integral rigorously transforms the continuous matrices into their discrete counterparts $\mathbf{\bar{A}}$ and $\mathbf{\bar{B}}$:
\begin{equation}
    \mathbf{\bar{A}} = \exp(\mathbf{\Delta} \mathbf{A}),
\end{equation}
\begin{equation}
    \mathbf{\bar{B}} = (\mathbf{\Delta} \mathbf{A})^{-1} (\exp(\mathbf{\Delta} \mathbf{A}) - \mathbf{I}) \cdot \mathbf{\Delta} \mathbf{B}.
\end{equation}
Following this discretization, the continuous differential system collapses into an elegant, discrete linear recurrence relation:
\begin{equation}
    h_k = \mathbf{\bar{A}}h_{k-1} + \mathbf{\bar{B}}x_k,
\end{equation}
\begin{equation}
    y_k = \mathbf{C}h_k.
\end{equation}

\subsection{The Selective Mechanism and Linear Complexity}
In standard State Space Models, matrices $\mathbf{\bar{A}}$, $\mathbf{\bar{B}}$, and $\mathbf{C}$ are invariant across time, which heavily limits their ability to process dynamic, context-dependent inputs. The profound architectural innovation of the Selective State Space Model is making the step size $\mathbf{\Delta}$, as well as matrices $\mathbf{B}$ and $\mathbf{C}$, strictly data-dependent functions of the input $x_k$. 

This "Selective Scan" mechanism endows the network with the mathematical capability to dynamically filter information: it can actively selectively remember critical structural features (e.g., the sharp boundary of the myocardium) and actively forget irrelevant artifacts (e.g., MRI bias field noise). Because Equation 5.13 is a pure recurrence relation that updates the hidden state $h_k$ without needing to compute the $L \times L$ pairwise interactions demanded by Vision Transformers, the computational complexity drops from quadratic $\mathcal{O}(L^2)$ to strictly linear $\mathcal{O}(L)$. 

The algorithmic implementation of this process is detailed in Algorithm \ref{alg:mamba}. 

\begin{algorithm}[hbt!]
\caption{Selective Scan Mechanism for 3D Cardiac Tensors}
\label{alg:mamba}
\begin{algorithmic}[1]
\REQUIRE Input serialized tensor $\mathbf{x} \in \mathbb{R}^{B \times L \times C}$, Discretization parameter $\mathbf{\Delta}$
\ENSURE Transformed spatio-temporal features $\mathbf{y} \in \mathbb{R}^{B \times L \times C}$
\STATE \textbf{Step 1: Data-Dependent Projection}
\STATE $\mathbf{B}, \mathbf{C} \leftarrow \text{Linear}(\mathbf{x})$, $\mathbf{\Delta} \leftarrow \text{Softplus}(\text{Linear}(\mathbf{x}))$
\STATE \textbf{Step 2: Discretization via ZOH}
\STATE $\mathbf{\bar{A}} = \exp(\mathbf{\Delta} \mathbf{A})$
\STATE $\mathbf{\bar{B}} = (\mathbf{\Delta} \mathbf{A})^{-1} (\exp(\mathbf{\Delta} \mathbf{A}) - \mathbf{I}) \cdot \mathbf{\Delta} \mathbf{B}$
\STATE \textbf{Step 3: Recurrent Parallel Scan}
\FOR{$t = 1$ \TO $L$}
    \STATE $h_t = \mathbf{\bar{A}}_t h_{t-1} + \mathbf{\bar{B}}_t x_t$
    \STATE $y_t = \mathbf{C}_t h_t$
\ENDFOR
\RETURN $\mathbf{y}$
\end{algorithmic}
\end{algorithm}

\section{Stochastic Data Augmentation and Spatial Regularization}
Deep learning models, particularly massive architectures encompassing millions of parameters, are inherently prone to memorizing the training cohort rather than learning generalizable, abstract features. Given the constrained scale of the ACDC dataset (100 patients), empirical risk minimization dictates the necessity of artificially expanding the topological diversity of the training distribution. This was achieved through a rigorous, stochastically driven 3D data augmentation pipeline utilizing the MONAI framework.

\subsection{Mathematical Formulation of 3D Affine Transformations}
To ensure the Nano-Mamba U-Net achieves rotational and scale invariance against varying patient orientations inside the MRI scanner, dynamic 3D affine transformations were applied during the data loading sequence. 

Let a continuous spatial coordinate in the 3D MRI volume be represented as a homogeneous column vector $\mathbf{v} = [x, y, z, 1]^T$. A spatial augmentation is mathematically defined by a $4 \times 4$ transformation matrix $\mathcal{M}_{affine}$. For stochastic spatial rotation, the network randomly samples angles $\theta_x, \theta_y, \theta_z \in [-15^\circ, 15^\circ]$ from a uniform distribution. The composite 3D rotation matrix $\mathbf{R}$ is the cross-product of the axial rotation matrices:
\begin{equation}
    \mathbf{R}_x(\theta_x) = 
    \begin{bmatrix}
    1 & 0 & 0 & 0 \\
    0 & \cos\theta_x & -\sin\theta_x & 0 \\
    0 & \sin\theta_x & \cos\theta_x & 0 \\
    0 & 0 & 0 & 1
    \end{bmatrix}, \quad
    \mathbf{R}_y(\theta_y) = 
    \begin{bmatrix}
    \cos\theta_y & 0 & \sin\theta_y & 0 \\
    0 & 1 & 0 & 0 \\
    -\sin\theta_y & 0 & \cos\theta_y & 0 \\
    0 & 0 & 0 & 1
    \end{bmatrix},
\end{equation}
\begin{equation}
    \mathbf{R}_z(\theta_z) = 
    \begin{bmatrix}
    \cos\theta_z & -\sin\theta_z & 0 & 0 \\
    \sin\theta_z & \cos\theta_z & 0 & 0 \\
    0 & 0 & 1 & 0 \\
    0 & 0 & 0 & 1
    \end{bmatrix}.
\end{equation}
Similarly, for stochastic scaling to simulate volumetric variations in dilated or hypertrophic hearts, a scale factor $s$ is uniformly sampled $s \sim \mathcal{U}(0.9, 1.1)$, defining the scaling matrix $\mathbf{S}$:
\begin{equation}
    \mathbf{S}(s) = \text{diag}(s, s, s, 1).
\end{equation}
The final augmented spatial coordinate $\mathbf{v}'$ is derived via sequential matrix multiplication:
\begin{equation}
    \mathbf{v}' = (\mathbf{S} \cdot \mathbf{R}_z \cdot \mathbf{R}_y \cdot \mathbf{R}_x) \mathbf{v}.
\end{equation}
To map the discrete voxel intensities to these newly transformed continuous coordinates without inducing topological tearing, trilinear interpolation was utilized for the image tensors, and nearest-neighbor interpolation was strictly maintained for the segmentation masks to preserve categorical integrity.

\subsection{Photometric Perturbations and Gaussian Noise}
Beyond geometric variations, the MRI sequences exhibit significant photometric inconsistencies due to fluctuating radio-frequency (RF) coil sensitivities and patient-specific magnetic susceptibilities. To regularize the network against these textural variations, stochastic Gaussian noise was injected directly into the normalized feature tensors.

Let $\mathcal{I}_{norm}$ represent the normalized MRI tensor. A noise matrix $\mathcal{N}$ of identical dimensions is sampled from a Gaussian probability density function (PDF):
\begin{equation}
    \mathcal{P}(x) = \frac{1}{\sigma \sqrt{2\pi}} \exp\left( -\frac{(x - \mu)^2}{2\sigma^2} \right),
\end{equation}
where $\mu = 0$ is the mean and $\sigma$ is the standard deviation, stochastically sampled between $[0.01, 0.1]$. The final photometrically augmented tensor is obtained via element-wise addition:
\begin{equation}
    \mathcal{I}_{aug} = \mathcal{I}_{norm} \oplus \mathcal{N}(\mu, \sigma).
\end{equation}
This rigorous augmentation policy forces the Nano-Mamba U-Net to rely on the overarching topological structure of the cardiac anatomy rather than overfitting to high-frequency, localized noise artifacts.

\section{Sliding Window Inference and Gaussian Blending}
While the hardware environment (RTX 4060, 8GB VRAM) was capable of handling small mini-batches during training, full-resolution volumetric inference on massive clinical MRI sequences often exceeds memory limits. To resolve this, a sliding-window inference topology was programmatically implemented during the evaluation phase.

The sliding-window algorithm decomposes the continuous macroscopic MRI volume $\mathcal{V}$ into overlapping microscopic 3D patches $\mathcal{P}_i$. Let $O \in (0, 1)$ denote the overlap ratio between adjacent patches (typically set to 0.5 to ensure 50\% spatial overlap). The Nano-Mamba architecture independently generates a semantic prediction mask $\hat{\mathcal{P}}_i$ for each patch.

However, simply concatenating these predicted patches results in severe boundary artifacts (checkerboard effects) due to the lack of global contextual awareness at the patch extremities. To mathematically synthesize a seamless, globally coherent volume, a Gaussian blending paradigm is utilized. For a predicted patch $\hat{\mathcal{P}}_i$ centered at coordinates $(c_x, c_y, c_z)$, an explicit Gaussian weighting matrix $\mathcal{W}$ of identical dimensions is generated:
\begin{equation}
    \mathcal{W}(x, y, z) = \exp \left( - \left( \frac{(x-c_x)^2}{2\sigma_x^2} + \frac{(y-c_y)^2}{2\sigma_y^2} + \frac{(z-c_z)^2}{2\sigma_z^2} \right) \right).
\end{equation}
This matrix assigns the maximum confidence weight (approaching 1.0) to the center voxels of the patch and exponentially decays the confidence weight (approaching 0.0) toward the patch boundaries. The final continuous volumetric prediction $\mathcal{V}_{pred}$ at any specific spatial coordinate $(x,y,z)$ is computed as the mathematically weighted average of all overlapping patches containing that coordinate:
\begin{equation}
    \mathcal{V}_{pred}(x,y,z) = \frac{\sum_{i} \hat{\mathcal{P}}_i(x,y,z) \cdot \mathcal{W}_i(x,y,z)}{\sum_{i} \mathcal{W}_i(x,y,z)}.
\end{equation}
This algorithmic synthesis ensures that the State Space representation retains extreme topological continuity across the entire cardiac cycle, eliminating localized edge artifacts and contributing directly to the superior Hausdorff Distance ($HD_{95}$) metrics achieved by the architecture.

\section{Formulation of the Objective Function}
Medical image segmentation inherently poses a severe class imbalance problem \cite{jadon2020survey}. Within a volumetric cardiac MRI tensor, the anatomical structures of interest (RV, MYO, LV) occupy a disproportionately minuscule fraction of the spatial domain, while the irrelevant background tissue heavily dominates the voxel count. Training a neural network exclusively on standard pixel-wise metrics often leads to highly skewed decision boundaries, where the model converges to a naive local minimum by predicting everything as background.

To mathematically counter this imbalance, the Nano-Mamba U-Net is optimized using a composite objective function, unifying the strengths of distribution-based and overlap-based loss metrics. The total loss $\mathcal{L}_{Total}$ is formulated as a linear amalgamation of the standard Cross-Entropy loss ($\mathcal{L}_{CE}$) and the soft Dice overlap loss ($\mathcal{L}_{Dice}$):
\begin{equation}
    \mathcal{L}_{Total} = \lambda_1 \mathcal{L}_{CE} + \lambda_2 \mathcal{L}_{Dice},
\end{equation}
with the empirical weighting coefficients standardized to $\lambda_1 = 1.0$ and $\lambda_2 = 1.0$.

The Cross-Entropy component independently evaluates the voxel-level probability distribution, providing stable gradients during the initial optimization phase:
\begin{equation}
    \mathcal{L}_{CE} = - \frac{1}{N} \sum_{i=1}^{N} \sum_{c=1}^{C} g_{ic} \log(p_{ic}),
\end{equation}
where $N$ is the voxel total, $C$ represents the four categorical classes, and $g_{ic} \in \{0, 1\}$ is the one-hot encoded clinical ground truth.

Conversely, the Dice Loss \cite{sudre2017generalised, milletari2016vnet} directly optimizes for structural overlap, fundamentally ignoring the absolute size of the background:
\begin{equation}
    \mathcal{L}_{Dice} = 1 - \frac{2 \sum_{i=1}^{N} \sum_{c=1}^{C} p_{ic} g_{ic} + \epsilon}{\sum_{i=1}^{N} \sum_{c=1}^{C} p_{ic}^2 + \sum_{i=1}^{N} \sum_{c=1}^{C} g_{ic}^2 + \epsilon}.
\end{equation}
The accumulated gradients derived from $\mathcal{L}_{Total}$ are backpropagated utilizing the Adam stochastic optimizer \cite{kingma2014adam}, which inherently adapts the learning rate for individual parameters using first and second moment estimates, ensuring robust convergence.

\section{Comprehensive Quantitative Evaluation Framework}
While the objective function directly optimizes the neural network's parameters during training, evaluating the clinical viability of the finalized model requires a multidimensional metric framework. Relying solely on volumetric overlap metrics is fundamentally insufficient for medical image segmentation, as they exhibit severe insensitivity to critical boundary outliers \cite{taha2015metrics}. Consequently, this research constructs a dual-faceted evaluation framework comprising spatial overlap metrics and geometric surface distance metrics.

\subsection{Volumetric Overlap and Statistical Metrics}
The primary metric for assessing region-based intersection is the Sørensen–Dice Similarity Coefficient (DSC). For binary evaluation of a specific class, let $TP$ denote True Positives, $FP$ False Positives, and $FN$ False Negatives:
\begin{equation}
    DSC = \frac{2 \cdot TP}{2 \cdot TP + FP + FN}.
\end{equation}
To provide a more granular statistical analysis of the model's diagnostic behavior, Sensitivity (Recall) and Specificity are evaluated. Sensitivity measures the proportion of actual anatomical tissue correctly identified:
\begin{equation}
    Sensitivity = \frac{TP}{TP + FN}.
\end{equation}
Conversely, Specificity evaluates the network's ability to correctly identify the background, heavily penalizing anatomical hallucinations (over-segmentation) in irrelevant tissue regions:
\begin{equation}
    Specificity = \frac{TN}{TN + FP}.
\end{equation}

\subsection{Geometric Surface Distance Metrics}
Clinical diagnosis, particularly the computation of myocardial mass and ventricular ejection fractions, relies heavily on the precise localization of anatomical boundaries. A model may achieve a high DSC while predicting an aberrant, jagged boundary. 

Let $S_P$ represent the set of surface points on the predicted segmentation mask, and $S_G$ represent the set of surface points on the clinical ground truth. The shortest Euclidean distance from a single point $p \in S_P$ to the ground truth surface $S_G$ is formulated as:
\begin{equation}
    d(p, S_G) = \min_{g \in S_G} \| p - g \|_2.
\end{equation}
The Average Symmetric Surface Distance (ASSD) computes the mean of all minimal distances from $S_P$ to $S_G$ and vice versa, providing a robust measure of overall contour alignment:
\begin{equation}
    ASSD(S_P, S_G) = \frac{1}{|S_P| + |S_G|} \left( \sum_{p \in S_P} d(p, S_G) + \sum_{g \in S_G} d(g, S_P) \right).
\end{equation}

Furthermore, to evaluate the worst-case boundary errors (e.g., a critical single-voxel leakage into an adjacent organ), the Hausdorff Distance (HD) is utilized. HD measures the maximum of all minimal distances between the two surfaces. To prevent extreme susceptibility to single-voxel noise artifacts inherent in MRI, the 95th percentile of the Hausdorff Distance ($HD_{95}$) is employed as the standard robustness metric:
\begin{equation}
    HD(S_P, S_G) = \max \left\{ \sup_{p \in S_P} d(p, S_G), \sup_{g \in S_G} d(g, S_P) \right\}.
\end{equation}
A lower $HD_{95}$ and $ASSD$ mathematically guarantee that the predicted boundaries are tightly conforming to the expert delineations.

\section{Computational Environment and Accelerated Ecosystem}
The empirical execution of state space models and massive 3D volumetric tensors necessitates a highly optimized software-hardware synergy. This research explicitly outlines the computational ecosystem designed to facilitate the rapid prototyping and edge-deployment testing of the Nano-Mamba U-Net.

\subsection{Software Infrastructure and Hardware Constraints}
The architectural implementation and automatic differentiation were facilitated by the PyTorch deep learning framework, utilizing its dynamic Directed Acyclic Graph (DAG) for the Autograd engine, which is critical for compiling the recurrent Mamba operations. 

All training, validation, and computational benchmarking (e.g., FPS profiling) were strictly constrained to a single consumer-grade NVIDIA GeForce RTX 4060 GPU, equipped with 8 Gigabytes of GDDR6 Video RAM (VRAM). This deliberate hardware constraint aligns with the thesis objective of proving the feasibility of advanced SOTA cardiac diagnostic AI on portable, low-power edge clinical devices.

\subsection{Mixed Precision Training and Hyperparameter Configuration}
Processing 3D continuous MRI tensors naturally imposes a severe memory bottleneck. Standard full-precision (FP32) tensors for large batch sizes quickly induce Out-Of-Memory (OOM) failures on an 8GB VRAM architecture. To circumvent this, the training protocol employed Automatic Mixed Precision (AMP) \cite{micikevicius2017mixed}. AMP dynamically casts specific tensor operations (such as linear layer multiplications and 3D convolutions) from 32-bit floating-point (FP32) down to 16-bit half-precision (FP16), while strictly preserving critical accumulation layers and gradient updates in FP32 to prevent numerical underflow.
\begin{equation}
    Memory(FP16) \approx \frac{1}{2} Memory(FP32).
\end{equation}

To ensure the reproducibility of the Nano-Mamba U-Net's 95.83\% DSC performance, the training environment and hyperparameter configurations were strictly standardized. All weights were initialized using the Kaiming Uniform distribution to maintain a stable variance of activations across layers. The global hyperparameter configuration is detailed in Table \ref{tab:hyperparams}.

\begin{table}[hbt!]
\caption{Global Hyperparameter Configuration for Nano-Mamba U-Net Training.}
\label{tab:hyperparams}
\centering
\begin{tabular}{l | p{8cm}}
\hline
\textbf{Hyperparameter} & \textbf{Selection Rationale and Value} \\
\hline
Initial Learning Rate & $1 \times 10^{-4}$ (Selected to ensure stable convergence in the early epochs). \\
Batch Size & 1 (Constrained by the volumetric 3D tensor dimensions and 8GB VRAM limit). \\
Optimizer & Adam ($\beta_1=0.9, \beta_2=0.999$, $\epsilon=1e-8$). \\
Weight Decay & $1 \times 10^{-5}$ (Implemented to prevent spatial overfitting in the 100-patient cohort). \\
Epochs & 150 (Total iterations required for loss stabilization). \\
Data Augmentation & Random 3D Rotation ($\pm 15^\circ$), Scaling ($[0.9, 1.1]$), and Gaussian Noise. \\
\hline
\end{tabular}
\end{table}

The choice of a mini-batch size of 1 is a strategic clinical constraint. In point-of-care environments, data is often processed patient-by-patient; thus, optimizing the model for singular volumetric inference ensures its immediate readiness for edge-device deployment.
